%% kapa: exam=1 round=33 year=2022 period=1 form=A subject=economics subject_ko=경제학원론 seq=6 number=46
%% source: Q-Net 감정평가사 기출문제 (cst00309) artlSeq=5209997 file=raw/1008/5209997/2022년도 제33회 감정평가사 1차시험 1교시(민법,경제학,부동산학).pdf
%% conversion: kordoc 4.24.51 -> tools/gen1.py (auto); review=proofread-sonnet; flags=tikz
\begin{question}{46}
완전경쟁시장에서 이윤극대화를 추구하는 기업 A의 한계비용(MC), 평균총비용(AC), 평균가변비용(AVC)은 아래 그림과 같다. 시장가격이 P$_1$, P$_2$, P$_3$, P$_4$, P$_5$로 주어질 때, 이에 관한 설명으로 옳지 않은 것은?
\begin{center}
\begin{tikzpicture}[x=0.075mm,y=-0.075mm,font=\footnotesize,line width=0.7pt]
\draw[-{Stealth[length=1.6mm]}] (115,745) -- (115,45);
\draw[-{Stealth[length=1.6mm]}] (115,745) -- (975,745);
\foreach \y/\l in {353/1,487/2,548/3,584/4,636/5}
  \draw[dashed,line width=0.5pt] (115,\y) -- (995,\y) (115,\y) node[left=3pt] {P$_\l$};
\draw plot[smooth] coordinates {(172,208) (200,280) (240,355) (300,420) (370,460) (450,482) (550,485) (640,478) (720,445) (790,395) (840,340) (880,260) (908,150)};
\draw plot[smooth] coordinates {(130,430) (160,480) (220,535) (300,572) (400,583) (490,583) (580,568) (670,540) (750,488) (800,440) (838,375)};
\draw plot[smooth] coordinates {(148,660) (200,680) (275,693) (340,688) (400,665) (450,625) (500,568) (550,490) (580,420) (600,340) (607,255)};
\node at (912,122) {AC};
\node at (887,393) {AVC};
\node at (635,215) {MC};
\node at (70,45) {P};
\node at (1000,770) {Q};
\node at (75,768) {0};
\end{tikzpicture}
\end{center}
\begin{choices}
  \item P$_1$일 때 총수입이 총비용보다 크다.
  \item P$_2$일 때 손익분기점에 있다.
  \item P$_3$일 때 총수입으로 가변비용을 모두 충당하고 있다.
  \item P$_4$일 때 총수입으로 고정비용을 모두 충당하고 있다.
  \item P$_5$일 때 조업중단을 한다.
\end{choices}
\end{question}
\answer{4}
